\section{Dimension of fibres}
\label{section-dimension-fibres}

\noindent
We study the behaviour of dimensions of fibres, using
Zariski's main theorem. Recall that we defined the
dimension $\dim_x(X)$ of a topological space $X$ at a point $x$
in Topology, Definition \ref{topology-definition-Krull}.

\begin{definition}
\label{definition-relative-dimension}
Suppose that $R \to S$ is of finite type, and let
$\mathfrak q \subset S$ be a prime lying over a prime
$\mathfrak p$ of $R$.
We define the {\it relative dimension
of $S/R$ at $\mathfrak q$}, denoted
$\dim_{\mathfrak q}(S/R)$, to be the dimension
of $\Spec(S \otimes_R \kappa(\mathfrak p))$
at the point corresponding to $\mathfrak q$. We let
$\dim(S/R)$ be the supremum of $\dim_{\mathfrak q}(S/R)$
over all $\mathfrak q$. This is called the
{\it relative dimension of} $S/R$.
\end{definition}

\noindent
In particular, $R \to S$ is quasi-finite at $\mathfrak q$ if
and only if $\dim_{\mathfrak q}(S/R) = 0$. The following lemma
is more or less a reformulation of Zariski's Main Theorem.

\begin{lemma}
\label{lemma-quasi-finite-over-polynomial-algebra}
Let $R \to S$ be a finite type ring map.
Let $\mathfrak q \subset S$ be a prime.
Suppose that $\dim_{\mathfrak q}(S/R) = n$.
There exists a $g \in S$, $g \not\in \mathfrak q$
such that $S_g$ is quasi-finite over a
polynomial algebra $R[t_1, \ldots, t_n]$.
\end{lemma}

\begin{proof}
The ring $\overline{S} = S \otimes_R \kappa(\mathfrak p)$ is
of finite type over $\kappa(\mathfrak p)$.
Let $\overline{\mathfrak q}$ be the prime of $\overline{S}$
corresponding to $\mathfrak q$.
By definition of
the dimension of a topological space at a point there exists
an open $U \subset \Spec(\overline{S})$ with
$\overline{q} \in U$ and $\dim(U) = n$.
Since the topology on $\Spec(\overline{S})$ is
induced from the topology on $\Spec(S)$ (see
Remark \ref{remark-fundamental-diagram}), we can find
a $g \in S$, $g \not \in \mathfrak q$ with image
$\overline{g} \in \overline{S}$ such that
$D(\overline{g}) \subset U$.
Thus after replacing $S$ by $S_g$ we see that
$\dim(\overline{S}) = n$.

\medskip\noindent
Next, choose generators $x_1, \ldots, x_N$ for $S$ as an $R$-algebra. By
Lemma \ref{lemma-Noether-normalization}
there exist elements $y_1, \ldots, y_n$ in the $\mathbf{Z}$-subalgebra of $S$
generated by $x_1, \ldots, x_N$ such that the map
$R[t_1, \ldots, t_n] \to S$, $t_i \mapsto y_i$ has the property
that $\kappa(\mathfrak p)[t_1\ldots, t_n] \to \overline{S}$
is finite. In particular, $S$ is quasi-finite over $R[t_1, \ldots, t_n]$
at $\mathfrak q$. Hence, by Lemma \ref{lemma-quasi-finite-open}
we may replace $S$ by $S_g$ for some $g\in S$, $g \not \in \mathfrak q$
such that $R[t_1, \ldots, t_n] \to S$ is quasi-finite.
\end{proof}

\begin{lemma}
\label{lemma-refined-quasi-finite-over-polynomial-algebra}
Let $R \to S$ be a ring map. Let $\mathfrak q \subset S$
be a prime lying over the prime $\mathfrak p$ of $R$.
Assume
\begin{enumerate}
\item $R \to S$ is of finite type,
\item $\dim_{\mathfrak q}(S/R) = n$, and
\item $\text{trdeg}_{\kappa(\mathfrak p)}\kappa(\mathfrak q) = r$.
\end{enumerate}
Then there exist $f \in R$, $f \not \in \mathfrak p$,
$g \in S$, $g \not\in \mathfrak q$ and a quasi-finite ring map
$$
\varphi : R_f[x_1, \ldots, x_n] \longrightarrow S_g
$$
such that $\varphi^{-1}(\mathfrak qS_g) =
(\mathfrak p, x_{r + 1}, \ldots, x_n)R_f[x_1, \ldots, x_n]$
\end{lemma}

\begin{proof}
After replacing $S$ by a principal localization we may assume there
exists a quasi-finite ring map $\varphi : R[t_1, \ldots, t_n] \to S$, see
Lemma \ref{lemma-quasi-finite-over-polynomial-algebra}.
Set $\mathfrak q' = \varphi^{-1}(\mathfrak q)$.
Let $\overline{\mathfrak q}' \subset \kappa(\mathfrak p)[t_1, \ldots, t_n]$
be the prime corresponding to $\mathfrak q'$. By
Lemma \ref{lemma-refined-Noether-normalization}
there exists a finite ring map
$\kappa(\mathfrak p)[x_1, \ldots, x_n] \to
\kappa(\mathfrak p)[t_1, \ldots, t_n]$
such that the inverse image of $\overline{\mathfrak q}'$ is
$(x_{r + 1}, \ldots, x_n)$. Let
$\overline{h}_i \in \kappa(\mathfrak p)[t_1, \ldots, t_n]$
be the image of $x_i$. We can find an element
$f \in R$, $f \not \in \mathfrak p$
and $h_i \in R_f[t_1, \ldots, t_n]$ which map to $\overline{h}_i$
in $\kappa(\mathfrak p)[t_1, \ldots, t_n]$. Then the ring map
$$
R_f[x_1, \ldots, x_n] \longrightarrow R_f[t_1, \ldots, t_n]
$$
becomes finite after tensoring with $\kappa(\mathfrak p)$.
In particular, $R_f[t_1, \ldots, t_n]$ is quasi-finite over
$R_f[x_1, \ldots, x_n]$ at the prime $\mathfrak q'R_f[t_1, \ldots, t_n]$.
Hence, by
Lemma \ref{lemma-quasi-finite-open}
there exists a $g \in R_f[t_1, \ldots, t_n]$,
$g \not \in \mathfrak q'R_f[t_1, \ldots, t_n]$
such that $R_f[x_1, \ldots, x_n] \to R_f[t_1, \ldots, t_n, 1/g]$
is quasi-finite. Thus we see that the composition
$$
R_f[x_1, \ldots, x_n] \longrightarrow
R_f[t_1, \ldots, t_n, 1/g] \longrightarrow S_{\varphi(g)}
$$
is quasi-finite and we win.
\end{proof}

\begin{lemma}
\label{lemma-dimension-inequality-quasi-finite}
Let $R \to S$ be a finite type ring map.
Let $\mathfrak q \subset S$ be a prime lying over $\mathfrak p \subset R$.
If $R \to S$ is quasi-finite at $\mathfrak q$, then
$\dim(S_{\mathfrak q}) \leq \dim(R_{\mathfrak p})$.
\end{lemma}

\begin{proof}
If $R_{\mathfrak p}$ is Noetherian
(and hence $S_{\mathfrak q}$ Noetherian since it is essentially of
finite type over $R_{\mathfrak p}$)
then this follows immediately from
Lemma \ref{lemma-dimension-base-fibre-total} and the
definitions. In the general case, let $S'$ be the integral
closure of $R_\mathfrak p$ in $S_\mathfrak p$.
By Zariski's Main Theorem \ref{theorem-main-theorem}
we have $S_{\mathfrak q} = S'_{\mathfrak q'}$ for some
$\mathfrak q' \subset S'$ lying over $\mathfrak q$.
By Lemma \ref{lemma-integral-dim-up} we have
$\dim(S') \leq \dim(R_\mathfrak p)$ and hence a fortiori
$\dim(S_\mathfrak q) = \dim(S'_{\mathfrak q'}) \leq \dim(R_\mathfrak p)$.
\end{proof}

\begin{lemma}
\label{lemma-dimension-quasi-finite-over-polynomial-algebra}
\begin{slogan}
A quasi-finite cover of affine n-space has dimension at most n.
\end{slogan}
Let $k$ be a field. Let $S$ be a finite type $k$-algebra.
Suppose there is a quasi-finite $k$-algebra map
$k[t_1, \ldots, t_n] \subset S$. Then $\dim(S) \leq n$.
\end{lemma}

\begin{proof}
By Lemma \ref{lemma-dim-affine-space} the dimension of
any local ring of $k[t_1, \ldots, t_n]$ is at most $n$.
Thus the result follows from
Lemma \ref{lemma-dimension-inequality-quasi-finite}.
\end{proof}

\begin{lemma}
\label{lemma-dimension-fibres-bounded-open-upstairs}
Let $R \to S$ be a finite type ring map.
Let $\mathfrak q \subset S$ be a prime.
Suppose that $\dim_{\mathfrak q}(S/R) = n$.
There exists an open neighbourhood $V$ of $\mathfrak q$
in $\Spec(S)$ such that
$\dim_{\mathfrak q'}(S/R) \leq n$ for all $\mathfrak q' \in V$.
\end{lemma}

\begin{proof}
By Lemma \ref{lemma-quasi-finite-over-polynomial-algebra}
we see that we may assume that $S$ is quasi-finite over
a polynomial algebra $R[t_1, \ldots, t_n]$. Considering
the fibres, we reduce to
Lemma \ref{lemma-dimension-quasi-finite-over-polynomial-algebra}.
\end{proof}

\noindent
In other words, the lemma says that the set of points where the
fibre has dimension $\leq n$ is open in $\Spec(S)$.
The next lemma says that formation of this open commutes with
base change.
If the ring map is of finite presentation then this set is
quasi-compact open (see below).

\begin{lemma}
\label{lemma-dimension-fibres-bounded-open-upstairs-base-change}
Let $R \to S$ be a finite type ring map.
Let $R \to R'$ be any ring map.
Set $S' = R' \otimes_R S$ and denote $f : \Spec(S') \to \Spec(S)$
the associated map on spectra.
Let $n \geq 0$.
The inverse image
$f^{-1}(\{\mathfrak q \in \Spec(S) \mid
\dim_{\mathfrak q}(S/R) \leq n\})$
is equal to
$\{\mathfrak q' \in \Spec(S') \mid
\dim_{\mathfrak q'}(S'/R') \leq n\}$.
\end{lemma}

\begin{proof}
The condition is formulated in terms of dimensions
of fibre rings which are of finite type over a field.
Combined with
Lemma \ref{lemma-dimension-at-a-point-preserved-field-extension}
this yields the lemma.
\end{proof}

\begin{lemma}
\label{lemma-dimension-fibres-bounded-quasi-compact-open-upstairs}
Let $R \to S$ be a ring homomorphism of finite presentation.
Let $n \geq 0$. The set
$$
V_n = \{\mathfrak q \in \Spec(S) \mid \dim_{\mathfrak q}(S/R) \leq n\}
$$
is a quasi-compact open subset of $\Spec(S)$.
\end{lemma}

\begin{proof}
It is open by Lemma \ref{lemma-dimension-fibres-bounded-open-upstairs}.
Let $S = R[x_1, \ldots, x_n]/(f_1, \ldots, f_m)$ be a presentation of
$S$. Let $R_0$ be the $\mathbf{Z}$-subalgebra of $R$ generated by the
coefficients of the polynomials $f_i$.
Let $S_0 = R_0[x_1, \ldots, x_n]/(f_1, \ldots, f_m)$.
Then $S = R \otimes_{R_0} S_0$. By
Lemma \ref{lemma-dimension-fibres-bounded-open-upstairs-base-change}
$V_n$ is the inverse image of an open $V_{0, n}$ under the quasi-compact
continuous map $\Spec(S) \to \Spec(S_0)$. Since
$S_0$ is Noetherian we see that $V_{0, n}$ is quasi-compact.
\end{proof}

\begin{lemma}
\label{lemma-finite-type-domain-over-valuation-ring-dim-fibres}
Let $R$ be a valuation ring with residue field $k$ and field
of fractions $K$. Let $S$ be a domain containing $R$ such that
$S$ is of finite type over $R$. If $S \otimes_R k$ is not the
zero ring then
$$
\dim(S \otimes_R k) = \dim(S \otimes_R K)
$$
In fact, $\Spec(S \otimes_R k)$ is equidimensional.
\end{lemma}

\begin{proof}
It suffices to show that $\dim_{\mathfrak q}(S/k)$ is equal
to $\dim(S \otimes_R K)$ for every prime $\mathfrak q$ of
$S$ containing $\mathfrak m_RS$. Pick such a prime. By
Lemma \ref{lemma-dimension-fibres-bounded-open-upstairs}
the inequality $\dim_{\mathfrak q}(S/k) \geq \dim(S \otimes_R K)$
holds. Set $n = \dim_{\mathfrak q}(S/k)$. By
Lemma \ref{lemma-quasi-finite-over-polynomial-algebra}
after replacing $S$ by $S_g$ for some $g \in S$, $g \not \in \mathfrak q$
there exists a quasi-finite ring map
$R[t_1, \ldots, t_n] \to S$. If $\dim(S \otimes_R K) < n$,
then $K[t_1, \ldots, t_n] \to S \otimes_R K$ has a nonzero kernel.
Say $f = \sum a_I t_1^{i_1}\ldots t_n^{i_n}$. After dividing
$f$ by a nonzero coefficient of $f$ with minimal valuation, we may
assume $f\in R[t_1, \ldots, t_n]$ and some $a_I$ does not map
to zero in $k$. Hence the ring map $k[t_1, \ldots, t_n] \to S \otimes_R k$
has a nonzero kernel which implies that $\dim(S \otimes_R k) < n$.
Contradiction.
\end{proof}
























